% ============================================================
% 00-abstract.tex
% ============================================================
\begin{abstract}
University counseling combines relational curricula, cohort-specific fees, policy tables, and narrative regulations, for which a uniform evidence path may be unsuitable. We present CTU-Chat, a supervisor-routed multi-agent system that assigns queries to four specialists with scoped representations and tools: Neo4j for curricula, structured lookup and deterministic calculation for financial data, and hybrid document retrieval for regulations. A 1,960-decision semantic-neighbor stress test compares full-registry and Top-$k$ single-agent baselines with routed generic and specialist gates across registries of 11--51 tools. At 51 tools, the specialist exceeds the global Top-10 baseline in decision-level end-to-end pass rate by 9.3 percentage points (95\% CI $[+2.0,+17.3]$), but uses 70.8\% more input tokens and $2.03\times$ the latency. Relative to the full-registry baseline, it uses 52.4\% fewer input tokens, while its 4.7-point end-to-end difference has a CI containing zero. Positive high-minus-low interactions for selection and end-to-end pass rate also have CIs containing zero, so a scaling-rate advantage is not established. On a redesigned 100-query held-out set, the full retrieval stack achieves 0.7400 Hit@1, 0.8608 Recall@5, and 0.8333 MRR@10; the full generation configuration reaches 0.6534 Context Recall, 0.8636 Faithfulness, and 0.8325 Source Recall. These findings suggest that domain-specific disambiguation and heterogeneous evidence allocation can support institutional counseling under the evaluated conditions, with configuration-dependent accuracy and cost trade-offs.

\keywords{Multi-Agent Systems \and Domain-Specialized Agents \and Centralized Supervisor Routing \and Heterogeneous Knowledge Allocation \and Retrieval-Augmented Generation \and University Counseling}
\end{abstract}
